\section{Introduction and prior work}\label{sec:report-introduction}
\subsection{The original question: a proof of Proposition 28}
Bordenave's random-lift proof of Friedman's second-eigenvalue theorem
uses a bound on products of centered permutation entries
\cite[Proposition 28]{Bordenave}. The same probabilistic input appears in
Bordenave--Collins \cite[Proposition 22]{BC} and in the biregular
spectral-gap argument of Brito--Dumitriu--Harris \cite[Lemma 5.6]{BDH}.
The estimate and the graph eigenvalue theorems are established prior
work. Problem 21 asks for an intuitive proof of that estimate and an
explanation of its spectral role. Sections~\ref{sec:mixed} and
\ref{sec:spectral} give that answer before the new recovery results.

A permutation matrix has exactly one nonzero entry in each row and
column. Its entries have mean $1/N$, but they are dependent. For a
uniform permutation matrix $M$, two distinct diagonal cells satisfy
\begin{equation}\label{eq:two}
\E[(M_{11}-N^{-1})(M_{22}-N^{-1})]
=\frac{1}{N(N-1)}-\frac1{N^2}
=\frac{1}{N^2(N-1)}.
\end{equation}
Two cells in one row instead have centered expectation $-N^{-2}$.
The smaller diagonal scale is cancellation, not independence. Taking
absolute values before using the row and column constraints would lose it.

Our proof keeps isolated centered singletons together, conditions only
on a compatible set of prescribed matches, and sums a row of the
remaining smaller uniform permutation matrix. The resulting positive
recurrence permits the remaining factors to be expanded safely.
Theorem~\ref{thm:mixed} gives the exact mixed-product envelope;
Corollary~\ref{cor:mixedscalar} recovers the requested scalar estimate.
Theorems~\ref{thm:oneconflict} and \ref{thm:twoconflict} give attained
maxima with one or two conflicting singletons. The base recurrence,
parity-dependent leading terms, and normalized error recurrences are due
to Collins--Nagatsu \cite[Proposition 3.3]{CN}; we credit and reuse them.
The graph path enumeration and tangle estimates remain cited inputs,
rather than new graph-norm theorems proved here.

\subsection{Prior work on recovering hidden graph groups}
A separate line of work asks whether exact neighbor-count groups can be
found after their labels are erased. Algebraic fiber recovery is not a
new idea: Fellows--Koblitz discussed hidden fibers in $K_4$ covers
\cite{rec:KidKrypto}, and Bernstein--Cottaar--Lange used echelon
reduction and grouping of affine expressions to expose those fibers
\cite{rec:BCL}. Regular-block-model rigidity and spectral recovery
also have direct predecessors \cite{rec:BDGHT,rec:Barucca,rec:KP,BDH}.
Our result should be compared with these recovery results, rather than
with noiseless detection alone.

The counting tools likewise have established sources. Vontobel's
cover enumerators and entropy interpretation \cite{rec:Vontobel},
Mori--Tanaka's positive interior prefactors \cite{rec:MoriTanaka},
and M\'esz\'aros's finite-group mixing \cite{rec:Meszaros} are relevant
predecessors. The sparse linear algebra uses classical recurrence and
preconditioning methods \cite{rec:Wiedemann,rec:Chen}. We do not claim
to invent Gaussian elimination, cover entropy, finite-field moments,
or black-box linear algebra. Further comparisons for the original
moment refinements appear in Appendix~\ref{sec:moment-discussion}.

\subsection{Main new contribution: unknown-frame recovery}
The principal new result of this report is
Theorem~\ref{thm:report-recovery}: recover the unknown degree frame,
its number and sizes of classes, and its hidden equitable partition,
from the graph and fixed class-count/degree bounds alone. The graph is
sampled from independent uniform simple internal regular and cross
biregular factors. The sufficient certificate is
\[
 q_{ij}=q_{ji},\qquad 0\le q_{ij}<\tfrac12,\qquad
                          \sum_jd_{ij}q_{ij}=1.
\]
The weights certify the model but need not be supplied to the algorithm.
They may be zero. Unequal class sizes, unequal total degrees, and absent
internal factors are allowed. The theorem includes every fixed connected
simple regular template of degree at least three.

There are two different guarantees. With high probability, every equitable
partition with at most the supplied constant $K$ classes coarsens the
plant; hence the plant is the unique finest such partition. Independently
of this identifiability event, an always-correct algorithm lists every
valid bounded-class partition. Its unconditional expected cost is
$\widetilde O(n^2)$ bit operations, with expected
$\widetilde O(n)$ working bits. The proof counts unusual inputs with large
solution spaces, rather than deleting them before computing an average.
It requires neither a separated planted eigenvalue nor a one-dimensional
real eigenspace. This recovery theorem uses its own rigidity and raw
channel-moment arguments; it is not deduced from Proposition 28 alone.

\paragraph{Reading and proof route.}
First read Section~\ref{sec:mixed} for the complete Proposition 28 input
and Section~\ref{sec:spectral} for its spectral interpretation.
Section~\ref{sec:recovery-overview} then states the main recovery theorem
with examples, an algorithm explanation, and its meaning.
The complete recovery proofs immediately follow in
Sections~\ref{rec:sec:unknown-frame}--\ref{rec:sec:vector-algorithm}:
rigidity, bounded raw channel moments, certified full kernels, and the
expected-work composition. The later appendices collect quantitative
moment refinements and their proofs; they are supplementary to the main
recovery result.

\paragraph{Scope of the claims.}
All recovery parameters are fixed as the graph grows, the sampling laws
are specified and independent, and neighbor counts are exact. The
listing guarantee is always-correct; choosing the hidden plant also uses
the stated probabilistic structural event. Expected nearly quadratic
work is not a worst-case bound or a benchmark claim. Neither the
recovery theorem nor the moment estimates establish robustness to noisy
edges, growing-parameter uniformity, or a new second-eigenvalue bound.
Written proofs and internal checks support the stated claims; external
peer review and historical firstness are not certified.
